% GENERATED FILE. Edit canonical lessons, then run npm run generate:books.
% canonical-source-hash: fnv1a64:6527dfcfaa287486
% canonical-lessons: ES-C637-favorecer, ES-C637-perjudicar, ES-C637-rentable, ES-C637-conveniente, ES-C637-convenir

\chapter{A favor o en contra}
\label{ch:es-a-favor-o-en-contra}
\hlchaptermodality{637}

\begin{chapteropening}
\textbf{By the end of this chapter:} \emph{I can say what favours or harms someone and what suits them.}

Complete the last lesson of chapter 637: I can say what favours or harms someone and what suits them.
\end{chapteropening}

\section[favorecer]{favorecer — to favour; to suit}

\label{lesson:ES-C637-favorecer}



\begin{quote}
Before the new one: say the Spanish for honest, then the Spanish for to deserve.
\end{quote}

\subsection*{You'll want to know: favorecer}
\textbf{favorecer} — \textquotedblleft{}to favour; to suit\textquotedblright{}.

\begin{quote}
La nueva ley favorece a las familias.
\end{quote}

\begin{quote}
\emph{The new law favours families.}
\end{quote}

Like \emph{conocer}: \emph{favorezco}. \emph{Ese color te favorece} — that colour suits you.

\begin{cousinweb}
\emph{favorecer} is built on \emph{favor} with the ending \emph{-ecer}; \emph{favor} is a learned borrowing from Latin \emph{favor}, from \emph{favere} 'to be well disposed' (EtymDB).
\end{cousinweb}

\begin{grammarlens}[title={What you've built}]
One more word for telling, explaining and describing at B1.
\end{grammarlens}

\subsection*{Your turn}
\begin{itemize}
\raggedright
  \item \emph{Say it:} \emph{favorecer}
  \item \emph{Say it:} \emph{favorecer}, once more
  \item \emph{Say it:} a sentence of your own with \emph{favorecer}
  \item \emph{Recall:} say the Spanish for honest, then the Spanish for to deserve, then say \emph{favorecer} again
\end{itemize}

\begin{tcolorbox}[breakable,colback=teal!4,colframe=teal!35!black,title={Before you move on}]
What does favorecer mean? (\textquotedblleft{}To favour; to suit\textquotedblright{}.) Say it once more.
\end{tcolorbox}
\section[perjudicar]{perjudicar — to harm, to damage}

\label{lesson:ES-C637-perjudicar}



\begin{quote}
Before the new one: say the Spanish for to deserve, then the Spanish for to favour.
\end{quote}

\subsection*{You'll want to know: perjudicar}
\textbf{perjudicar} — \textquotedblleft{}to harm, to damage\textquotedblright{}.

\begin{quote}
Fumar perjudica la salud.
\end{quote}

\begin{quote}
\emph{Smoking harms your health.}
\end{quote}

Spelling: \emph{perjudiqué}. \emph{Perjudicial} — harmful.

\begin{cousinweb}
\emph{perjudicar} is a learned borrowing from Latin \emph{praeiudicare} 'to judge in advance, to harm a case', built on \emph{iudicare} and \emph{iudex} 'judge' (EtymDB).
\end{cousinweb}

\begin{grammarlens}[title={What you've built}]
One more word for telling, explaining and describing at B1.
\end{grammarlens}

\subsection*{Your turn}
\begin{itemize}
\raggedright
  \item \emph{Say it:} \emph{perjudicar}
  \item \emph{Say it:} \emph{perjudicar}, once more
  \item \emph{Say it:} a sentence of your own with \emph{perjudicar}
  \item \emph{Recall:} say the Spanish for to deserve, then the Spanish for to favour, then say \emph{perjudicar} again
\end{itemize}

\begin{tcolorbox}[breakable,colback=teal!4,colframe=teal!35!black,title={Before you move on}]
What does perjudicar mean? (\textquotedblleft{}To harm, to damage\textquotedblright{}.) Say it once more.
\end{tcolorbox}
\section[rentable]{rentable — profitable, worthwhile}

\label{lesson:ES-C637-rentable}



\begin{quote}
Before the new one: say the Spanish for to favour, then the Spanish for to harm.
\end{quote}

\subsection*{You'll want to know: rentable}
\textbf{rentable} — \textquotedblleft{}profitable, worthwhile\textquotedblright{}.

\begin{quote}
El negocio es rentable.
\end{quote}

\begin{quote}
\emph{The business is profitable.}
\end{quote}

Built on \emph{la renta} — income.

\begin{cousinweb}
\emph{rentable} is built on \emph{renta} 'income', which is usually traced to spoken Latin \emph{rendita} 'what is given back', from \emph{reddere} 'to give back' (Lewis \& Short).
\end{cousinweb}

\begin{grammarlens}[title={What you've built}]
One more word for telling, explaining and describing at B1.
\end{grammarlens}

\subsection*{Your turn}
\begin{itemize}
\raggedright
  \item \emph{Say it:} \emph{rentable}
  \item \emph{Say it:} \emph{rentable}, once more
  \item \emph{Say it:} a sentence of your own with \emph{rentable}
  \item \emph{Recall:} say the Spanish for to favour, then the Spanish for to harm, then say \emph{rentable} again
\end{itemize}

\begin{tcolorbox}[breakable,colback=teal!4,colframe=teal!35!black,title={Before you move on}]
What does rentable mean? (\textquotedblleft{}Profitable, worthwhile\textquotedblright{}.) Say it once more.
\end{tcolorbox}
\section[conveniente]{conveniente — advisable, convenient}

\label{lesson:ES-C637-conveniente}



\begin{quote}
Before the new one: say the Spanish for to harm, then the Spanish for profitable.
\end{quote}

\subsection*{You'll want to know: conveniente}
\textbf{conveniente} — \textquotedblleft{}advisable, convenient\textquotedblright{}.

\begin{quote}
Es conveniente llegar pronto.
\end{quote}

\begin{quote}
\emph{It is advisable to arrive early.}
\end{quote}

\emph{Es conveniente que} + subjunctive.

\begin{cousinweb}
\emph{conveniente} is a learned borrowing from Latin \emph{conveniens} 'agreeing, fitting', the participle of \emph{convenire} 'to come together', from \emph{venire} 'to come' (EtymDB).
\end{cousinweb}

\begin{grammarlens}[title={What you've built}]
One more word for telling, explaining and describing at B1.
\end{grammarlens}

\subsection*{Your turn}
\begin{itemize}
\raggedright
  \item \emph{Say it:} \emph{conveniente}
  \item \emph{Say it:} \emph{conveniente}, once more
  \item \emph{Say it:} a sentence of your own with \emph{conveniente}
  \item \emph{Recall:} say the Spanish for to harm, then the Spanish for profitable, then say \emph{conveniente} again
\end{itemize}

\begin{tcolorbox}[breakable,colback=teal!4,colframe=teal!35!black,title={Before you move on}]
What does conveniente mean? (\textquotedblleft{}Advisable, convenient\textquotedblright{}.) Say it once more.
\end{tcolorbox}
\section[convenir]{convenir — to suit, to be advisable; to agree}

\label{lesson:ES-C637-convenir}



\begin{quote}
Before the new one: say the Spanish for profitable, then the Spanish for advisable.
\end{quote}

\subsection*{You'll want to know: convenir}
\textbf{convenir} — \textquotedblleft{}to suit, to be advisable; to agree\textquotedblright{}.

\begin{quote}
Te conviene descansar.
\end{quote}

\begin{quote}
\emph{You would do well to rest.}
\end{quote}

Like \emph{venir}: \emph{conviene, convino}. \emph{Me conviene} — it suits me.

\begin{cousinweb}
\emph{convenir} is inherited from Latin \emph{convenire} 'to come together, to agree, to suit', built on \emph{venire} 'to come' (EtymDB).
\end{cousinweb}

\begin{grammarlens}[title={What you've built}]
One more word for telling, explaining and describing at B1.
\end{grammarlens}

\subsection*{Your turn}
\begin{itemize}
\raggedright
  \item \emph{Say it:} \emph{convenir}
  \item \emph{Say it:} \emph{convenir}, once more
  \item \emph{Say it:} a sentence of your own with \emph{convenir}
  \item \emph{Recall:} say the Spanish for profitable, then the Spanish for advisable, then say \emph{convenir} again
\end{itemize}

\begin{tcolorbox}[breakable,colback=teal!4,colframe=teal!35!black,title={Before you move on}]
What does convenir mean? (\textquotedblleft{}To suit, to be advisable; to agree\textquotedblright{}.) Say it once more.
\end{tcolorbox}
